
\subsection{Problem Statement}

The main objective of this project is to investigate whether students' academic efficiency depends solely on the amount of time dedicated to studying, or if it results from a combination of factors such as focus, quality of rest, digital habits, and overall well-being. This problem is particularly relevant in modern educational contexts, where students are exposed to multiple sources of distraction and varying lifestyle patterns that may impact their academic performance.

\subsection{Target Variables}

To address the defined problem, two complementary predictive scenarios are considered.

\textbf{Regression Scenario:}  
The variable \texttt{focus\_index} is used as the target for regression, representing a continuous measure of a student's ability to maintain concentration during study activities.

\textbf{Classification Scenario:}  
The variable \texttt{exam\_score} is transformed into a binary target representing academic performance. Specifically, a threshold is defined to distinguish between high performance (1) and normal or low performance (0), allowing the problem to be framed as a classification task.

\subsection{Stakeholders}

Several stakeholders can benefit from the outcomes of this study.

\textbf{Students:}  
Students can gain insights into whether increasing study time alone leads to better performance, or if other factors such as focus and well-being play a more significant role.

\textbf{Mentors and Tutors:}  
Educators and tutors can use the findings to recommend more effective and personalized study strategies, focusing not only on study duration but also on behavioral and lifestyle factors.

\textbf{Academic Support Services:}  
Institutional support services can leverage these insights to design interventions aimed at improving academic success, emphasizing healthy habits, balanced digital usage, and overall student well-being rather than solely increasing study workload.
